\section*{Erd\H{o}s Problem 405}

\paragraph{Statement and status.}
For a fixed odd prime \(p\), are there only finitely many positive integers
\(a\) for which \((p-1)!+a^{p-1}\) is a power of \(p\)?
The stronger classification is known: the solutions \((p,a,k)\) of
\((p-1)!+a^{p-1}=p^k\) are
\((3,1,1),(3,5,3),(5,1,2)\), by Yu and Liu.
The full classification and general finiteness are unresolved in this
attempt. Here are complete restrictions and a complete treatment of \(p=5\).

\paragraph{Prime-divisor and congruence restrictions.}
Let \(F=(p-1)!\). If a prime \(q<p\) divides \(a\), it divides
\(F+a^{p-1}=p^k\), which is impossible. If \(p\mid a\), Wilson's theorem
gives \(F+a^{p-1}\equiv-1\pmod p\), also impossible. Thus
\(\gcd(a,p!)=1\).

For \(p\geq5\), this makes \(a\) odd. Since \(8\mid F\) and
\(a^{p-1}\equiv1\pmod8\), we have \(p^k\equiv1\pmod8\).
Also \(3\mid F\), \(3\nmid a\), and \(p-1\) is even, so
\(p^k\equiv1\pmod3\). Consequently, if \(k\) is odd, then
\[
 p\equiv1\pmod{24}.
 \tag{1}
\]

\paragraph{All even exponents reduce to a finite search.}
If \(k=2h\), set \(X=p^h\), \(Y=a^{(p-1)/2}\). Then
\[
 (X-Y)(X+Y)=F,\qquad X>Y>0.
\]
Every solution comes from a positive factor pair \(u<v\), \(uv=F\),
with equal parity, by \(X=(u+v)/2\), \(Y=(v-u)/2\).
There are only finitely many such pairs, and
\(X\leq(F+1)/2\). In particular,
\[
 k\leq2\log_p((F+1)/2).
\]
This proves finiteness of the even-\(k\) portion for every fixed \(p\).

\paragraph{The complete case \(p=5\).}
Here \(a\) is odd and \(a^4+24\equiv9\pmod{16}\).
The powers of \(5\pmod{16}\) are \(5,9,13,1\), so \(k\equiv2\pmod4\).
Thus \(k\) is even, and the preceding factorization applies with
\(F=24\), \(Y=a^2\). The positive equal-parity factor pairs \(u<v\)
of \(24\) are \((2,12)\) and \((4,6)\). The first gives
\((X,Y)=(7,5)\), which fails both required power conditions.
The second gives \((X,Y)=(5,1)\), hence \(a=1,k=2\).
Indeed \(4!+1=25\).

\paragraph{Remaining gap and source.}
For \(p=3\), the two displayed solutions follow from
\(1^2+2=3\) and \(5^2+2=27\); their exhaustiveness is not proved here.
For primes \(p\equiv1\pmod{24}\), (1) leaves an unbounded odd-exponent
case untreated. The known full result is Kunrui Yu and Dehua Liu,
\emph{A complete resolution of a problem of Erd\H{o}s and Graham},
Rocky Mountain Journal of Mathematics (1996), 1235--1244,
\url{https://doi.org/10.1216/rmjm/1181072047}.
Statement checked at \url{https://www.erdosproblems.com/405}
on 24 September 2026.

\paragraph{Completion estimate.}
Even exponents and \(p=5\) are settled here; the general odd-exponent argument is missing.
\textbf{COMPLETION: 30\%}
